\section{Introduction}\label{sec:intro}

In audit, compliance, and clinical settings, a model that cannot show a
person \emph{why} it concluded something is often unusable regardless of
its accuracy. Regulators demand justifications for individual decisions,
auditors must defend every flag they raise, and the case for inherently
interpretable models over post-hoc explanations of black boxes in
high-stakes settings is well established~\cite{rudin2019}. At the same
time, the systems that offer strong probabilistic semantics for logical
knowledge --- probabilistic logic programming and statistical relational
learning (Section~\ref{sec:related}) --- present two obstacles to exactly
these users: their inference machinery (knowledge compilation, MRF
inference) is opaque to the domain practitioner who must sign off on a
decision, and their rule weights are fixed at inference time, so a rule's
reliability cannot respond to the operational context in which it fires.

\pla{} explores the opposite corner of the design space. It is a
rule-confidence engine simple enough to trace by hand --- a few hundred
lines of dependency-free Python --- that is honest about \emph{not} being
a probability distribution, and that invests the saved complexity in
three places: context-conditioned rule weights, learnability, and
measured explanation quality.

\paragraph{Contributions.}
\begin{enumerate}
\item \textbf{An explicit confidence calculus with verified properties}
  (Section~\ref{sec:calculus}). All operators are named and formally
  defined; the definitions are executable --- the test suite extracts and
  runs every formula in the semantics document --- and fixpoint
  convergence is machine-checked on 1{,}050 random cyclic rule systems in
  addition to a proof sketch.
\item \textbf{Context-conditioned, learnable rule weights}
  (Section~\ref{sec:learning}). Context variables contribute additive
  log-odds evidence, avoiding the saturation that afflicts
  multiply-and-cap certainty-factor updates; on binary facts the model
  class collapses to logistic regression over rule firings, so
  maximum-likelihood fitting is convex, provably monotone under
  full-batch gradient descent, and recovers planted ground truth.
\item \textbf{Measured explanation fidelity}
  (Section~\ref{sec:evaluation}). Rule traces are the engine's primary
  output, and their fidelity is quantified with deletion-based
  comprehensiveness and sufficiency against a reversed-attribution
  control, rather than asserted.
\item \textbf{A reproducibility discipline}
  (Section~\ref{sec:implementation}). Every number reported in this
  paper regenerates byte-for-byte from committed scripts; claims about
  competitor capabilities are themselves verified by tests against the
  installed competitor packages; and Appendix~\ref{app:provenance} maps
  each substantive claim to the artifact that verifies it.
\end{enumerate}

Section~\ref{sec:related} positions \pla{} against probabilistic logic
programming, weighted-logic learning, and its own certainty-factor
lineage. Sections~\ref{sec:calculus} and~\ref{sec:learning} present the
calculus and the learning theory. Section~\ref{sec:implementation}
describes the tool. Section~\ref{sec:evaluation} reports development
results on synthetic data --- including a negative finding we consider
informative --- and specifies the pending real-data study.
Section~\ref{sec:limitations} lists limitations explicitly.
